% 第三部分 C：摘要（由 main_zh.tex 引入）。
树状 rollout 为每个提示采样少数几个前缀并从中分支，各分支共享前缀的词元。一个能认出有希望的部分解答的评判模型，可以决定从哪里分支。我们在 MATH 上用 Qwen3.5-4B-Base 检验这一点。离线时，\kev{}、同样大小的自回归评判模型和策略自身的对数概率，在不同题目之间都能预测一个 128 词元的前缀最终是否正确（AUROC \NUM{auc_128_logprob}--\NUM{auc_128_artrack}），但在同一道题的前缀之间只能弱地区分（\NUM{aucw_128_kevcorrect}--\NUM{aucw_128_logprob}）；保留 \kev{} 在四个前缀中排得最高的两个，让正确续写的比例提高 \NUM{gainpt_128_kevtrack} 个百分点，策略的对数概率为 \NUM{gainpt_128_logprob}。在使用树状 rollout 的 GRPO 中（四个 128 词元的根，其中两个各续写四次，200 步，两个种子），\kev{} 选根让正确叶子的比例提高 \NUM{auditgain_all_treekev} 个百分点 \ci{\NUM{auditgainlo_all_treekev}}{\NUM{auditgainhi_all_treekev}}，随机选根为 \NUM{auditgain_all_treerand}；最终策略在五个基准中的四个上与随机选根一样准确（GSM8K：\NUM{dvtreerand_treekev_gsm} \ci{\NUM{dvtreerandlo_treekev_gsm}}{\NUM{dvtreerandhi_treekev_gsm}}）。按种子均值，两种树状条件最终都高于扁平 GRPO（MATH-500 上高 \NUM{d_treerand_math} 和 \NUM{d_treekev_math} 个百分点，AIME24 上高 \NUM{d_treerand_aime24} 和 \NUM{d_treekev_aime24}），但每种都有一个种子明显高于扁平采样、另一个与它持平，所以树形是否有帮助仍无定论。使用 \kev{} 选根的一个种子逐渐在答案之后重复解答，而验证器的奖励不惩罚这种重复。已撤回的早期版本的剪枝规则——评判模型的 $p(\text{yes})<0.3$ 时剪掉前缀——只剪掉 \NUM{thrpruned_128_kevtrack} 的前缀。本版取代了结果并非由实验产生的早期版本。代码、日志和本文的每一个数字都在仓库中。
